\section{Calibration and measurement}
\label{sec:companion-calibration}

Table~\ref{tab:companion-parameters} imports the reference values documented
in the quantitative section of \citet{pardo2026future}. It separates
macroeconomic benchmarks from illustrative AI parameters. Reusing these
values preserves comparability; it does not make them empirical estimates
of this autonomous-research economy.

\begin{table}[H]
\centering
\caption{Inherited parameters and their evidential status}
\label{tab:companion-parameters}
\small
\setlength{\tabcolsep}{4pt}
\renewcommand{\arraystretch}{1.12}
\begin{tabular}{P{0.10\textwidth}P{0.11\textwidth}P{0.26\textwidth}P{0.42\textwidth}}
\toprule
Parameter & Reference & Interpretation & Source or status\\
\midrule
$\alpha$ & 0.33 & Capital elasticity & Inherited one-third macroeconomic benchmark.\\
$\delta$ & 0.05 & Annual depreciation & Inherited benchmark; PWT is the comparison source, not a direct estimate of this value \citep{pwt110}.\\
$\rho$ & 0.04 & Annual discount rate & Preference calibration, not a measured market interest rate.\\
$n$ & 0.003 & Annual population growth & Inherited constant-rate approximation to the UN 2024--2100 world projection \citep{unwpp2024}.\\
$\gamma$ & 0.01 & Annual labor-augmenting progress & Illustrative trend, with the OECD long-run scenarios as a scale comparison \citep{oecdlongrun2025}.\\
$\omega_X$ & 0.20 & AI production weight & Illustrative, not estimated. Not the AI revenue share itself.\\
$\omega_L$ & 0.80 & Effective-labor weight & Defined by $1-\omega_X$.\\
$\eta$ & 0.20 & Research-services elasticity & Illustrative, not estimated.\\
$\sigma$ & 1 & AI--labor substitution & Maintained restriction of this paper.\\
$\chi$ & Free, $>0$ & Research productivity & Does not affect the BGP rates or expenditure ratios. Needs units and level targets for a quantitative level comparison.\\
\bottomrule
\end{tabular}
\end{table}

The population parameter extrapolates a finite-period approximation, not a
UN prediction of positive population growth forever. Likewise, measured
labor-productivity growth is not the same object as the growth of $A$ in
this model. The OECD comparison does not directly identify $\gamma$.
Both $n$ and $\gamma$ therefore belong in long-run sensitivity analysis.

The finite-frontier paper uses $\chi=1.4378$ to target an illustrative
transition duration. I do not transfer that timing target to the uncapped
economy. Equations~\eqref{eq:companion-gy}--\eqref{eq:companion-bgp-consumption}
are independent of $\chi$, whereas the initial BGP levels constructed in
Appendix~\ref{app:companion-proofs} depend on it. Calibrating $\chi$ cannot
select an arbitrary permanent growth rate.

At the inherited reference values, the analytical calculations give
\begin{equation}
 g_{Y/N}^*=g_w^*\simeq1.0867\%\ \text{per year},\qquad
 r^*\simeq5.0867\%\ \text{per year}.
 \label{eq:companion-reference-results}
\end{equation}
The growth increase relative to $\gamma=1\%$ is about $0.0867$ percentage
points per year. This comparison is between long-run model benchmarks, not
a transition from today's economy and not an estimate of AI's observed impact.

Empirical discipline requires a mapping from observables to model objects.
At unit elasticity, \eqref{eq:companion-distribution} implies that a measured
AI revenue share would map to $\omega_X$ only after dividing by $1-\alpha$,
provided the measurement matches the model's industry boundary and pricing
assumptions. AI investment, compute expenditure, and AI revenue are different
objects. Measurement difficulties discussed by \citet{oecdaiinvestment2025}
motivate caution; they do not select $\omega_X=0.20$.

For research, taking logs of \eqref{eq:companion-research-law} gives
\begin{equation}
 \log\dot B=\log\chi+\eta\log B+\eta\log M.
 \label{eq:companion-measurement-law}
\end{equation}
The left side is the log of an efficiency increment, not the growth rate
$g_B$. Moreover, research compute is an endogenous choice. Historical
algorithmic improvements cannot be read directly as an estimate of $\eta$:
they may also reflect human research, data, and complementary inputs omitted
from the autonomous benchmark. Identifying defensible ranges for $\eta$ and
$\omega_X$ remains a separate empirical task.
